\abstracttext{As audiências públicas da Câmara dos Deputados chegam ao público sobretudo por matérias
jornalísticas que selecionam poucas opiniões, sem indicar onde cada uma aparece na transcrição. Este
trabalho investiga se um perfil escrito só com as falas anteriores de uma pessoa ajuda um modelo de
linguagem a reconhecer o que a imprensa lhe atribuiu numa audiência posterior. A Unidade Deliberativa
Verificável localiza na fala da própria pessoa evidência candidata para 2.105 das 2.203 opiniões do
PublicHearingBR, com a localização na transcrição, o tipo e a origem do vínculo. Os perfis por ator,
escritos por um modelo exclusivamente com as falas de cada pessoa nas audiências de treino, têm
blocos de posições, critérios e valores, alinhamentos declarados e forma de argumentar. Na simulação,
o perfil e trechos recuperados das falas anteriores orientam um modelo a responder como a pessoa. Na
geração aberta, o modelo declara antes se o material trata do assunto e, depois da fala, indica o
item ou trecho que apoia cada frase, com o início de cada cópia conferido no texto de origem. A
avaliação usa como gabarito as opiniões publicadas sobre audiências posteriores, em perguntas de
múltipla escolha com distratores da mesma audiência. Nelas, a \textit{classifier-free guidance}
combina as probabilidades das quatro opções com e sem o material da pessoa e desfavorece as que o
modelo já escolheria só com o nome e o cargo.
<Resultados principais: acerto de cada condição no teste e diferenças entre condições com intervalo
de confiança.>}
